\begin{statementbox}
  \textbf{\textcolor{CStatement}{这个结论解决什么问题？}}

  一道函数值域题里，如果 $x$ 既出现在根号内，又出现在根号外，直接观察最高值、最低值往往很难。换元就是\textbf{把反复出现的一整块式子当成一个新字母}，把题目变成学过的二次函数或三角函数问题。

  \medskip
  \begin{center}\begin{tikzpicture}[font=\small,>=stealth,line join=round]
\draw[gray!45,rounded corners=3pt] (0,0) rectangle (14.2,9.4);
\node at (7.1,8.8) {整个根号作新变量：沿曲线读出对应范围};

\begin{scope}[shift={(2.3,2.5)},x=1cm,y=1.7cm]
\draw[->] (0,0)--(5.7,0) node[right]{$x$};\draw[->] (0,-.2)--(0,2.8) node[above]{$t$};
\draw[teal!80!black,thick,domain=0:2.15,samples=70] plot({\x*\x+1},\x);\draw[blue,very thick,domain=1:2,samples=50] plot({\x*\x+1},\x);
\foreach \x/\y in {2/1,5/2}{\draw[dashed] (\x,0)--(\x,\y)--(0,\y);\fill[blue] (\x,\y) circle(1.7pt);\node[below] at (\x,0){$\x$};\node[left] at (0,\y){$\y$};}
\node[below] at (1,0){$1$};\fill (1,0) circle(1.5pt);
\end{scope}
\node[align=left,anchor=west] at (8.5,5.1){$t=\sqrt{x-1}\ge0$\\[10pt] $x=t^2+1$\\[12pt] 全部：$x\ge1\iff t\ge0$\\[10pt] 限制：$2\le x\le5$\\[5pt] $\phantom{限制：}\iff1\le t\le2$};
\node[text width=13.2cm,align=center] at (7.1,.65) {粗蓝线表示额外限制后的部分；换元必须同时保留它的实际范围。};
\end{tikzpicture}
\end{center}
\end{statementbox}
\clearpage
\begin{statementbox}
\textbf{\textcolor{CStatement}{使用前先检查}}
  \begin{itemize}[leftmargin=*,itemsep=0.35em]
    \item 根号里的数必须大于或等于 $0$；分母不能为 $0$。
    \item 换元后，新字母能取哪些值，必须从原题条件推出来。
    \item 下文写 $a>0$；若题中给出 $a<0$，先把相关平方中的 $a$ 换成 $|a|$。
  \end{itemize}

  \medskip
  \textbf{\textcolor{CStatement}{结论一：根号里是一次式，直接设整个根号}}

  对于 $\sqrt{mx+n}$（$m\ne0$），令
  \[
    t=\sqrt{mx+n}\quad(t\ge0),\qquad x=\frac{t^2-n}{m}.
  \]
  用 $t$ 表示原式，再在 $t$ 的实际范围内求值域。\textbf{注意：若题目另有限制，$t\ge0$ 还不够。}

  \medskip
  \textbf{\textcolor{CStatement}{结论二：根号里是两个平方的差，像一个圆}}

  当 $a>0$、$-a\le x\le a$ 时，令
  \[
    x=a\sin\theta,\quad -\frac{\pi}{2}\le\theta\le\frac{\pi}{2},
    \qquad\sqrt{a^2-x^2}=a\cos\theta.
  \]
  于是 $x$ 与根号可以同时写成正弦、余弦，适合处理它们相加、相减的题。

  \medskip
  \textbf{\textcolor{CStatement}{结论三：根号里是两个平方的和}}

  当 $a>0$、$x$ 为任意实数时，可令
  \[
    x=a\tan\theta,\quad-\frac{\pi}{2}<\theta<\frac{\pi}{2},
    \qquad \sqrt{x^2+a^2}=\frac{a}{\cos\theta}.
  \]
  也可以不用三角函数：令 $t=\sqrt{x^2+a^2}+x>0$，则
  \[
    x=\frac{t^2-a^2}{2t},\qquad
    \sqrt{x^2+a^2}=\frac{t^2+a^2}{2t}.
  \]

  \medskip
  \textbf{\textcolor{CStatement}{结论四：根号里是 $x^2-a^2$，先分正负}}

  当 $a>0$ 时，必须有 $x\ge a$ 或 $x\le-a$。
  \begin{itemize}[leftmargin=*,itemsep=0.35em]
    \item 若 $x\ge a$，令 $t=x+\sqrt{x^2-a^2}\ge a$，则
      \[
        x=\frac{t^2+a^2}{2t},\qquad
        \sqrt{x^2-a^2}=\frac{t^2-a^2}{2t}.
      \]
    \item 若 $x\le-a$，令 $t=-x+\sqrt{x^2-a^2}\ge a$，则
      \[
        x=-\frac{t^2+a^2}{2t},\qquad
        \sqrt{x^2-a^2}=\frac{t^2-a^2}{2t}.
      \]
  \end{itemize}

  \textbf{\textcolor{CStatement}{最重要的一句话：}}\textbf{换元可以改写式子，但不能改掉原题允许的取值。}
\end{statementbox}
